%% toc: P3 38 38.1,38.2,38.3
\chapter{Connectedness and Compactness Are Not Expressible in the Topological Language}
\label{sec:p03-w03-p03-051-nondefinability}

The guarded language \(L_t\) of \cref{def:p03-topological-languages}
expresses the separation axioms, discreteness and the Alexandroff property
(\cref{thm:p03-topological-property-formulas}). This chapter proves that it
cannot express connectedness, and by the same method that it cannot express
compactness. Unrestricted \(L_2\) does express connectedness, so the result
marks a real loss of expressive power, paid for the basis invariance of
\cref{thm:p03-topological-basis}.

The inputs are a topological model in the sense of
\cref{def:p03-topological-carrier} whose space has a countable basis, and a
property of topological spaces. The conclusion is a second topological
model, on a countable subspace, that satisfies exactly the same
\(L_t\)-sentences and lacks the property. The tools are the basis
invariance theorem, the extent construction used in the proof of
\cref{thm:p03-topological-compactness-ls}, and ordinary downward
L\"owenheim--Skolem (\cref{thm:downward-ls}).

The statement for connectedness, and the definition of connectedness
used here, are those of the October~2006 draft of ten~Cate,
Gabelaia and Sustretov
\cite[Definition~13, p.~11; Theorem~2.9, p.~12]{TenCateGabelaiaSustretov2006TopologyDraft}.
The draft gives no proof; it attributes the theorem to the 1980 monograph
of Flum and Ziegler \cite{FlumZiegler1980TopologicalModelTheory}. The
proof below is constructed from the results already proved in
\cref{sec:p03-topological-model-theory}; it is not a reconstruction of the
monograph's argument, and the compactness corollary is not taken from either
work.

\section{Expressibility and standing assumptions}
\label{sec:p03-w03-p03-051-expressibility}

\begin{hbconvention}{def:p03-w03-p03-051-assumptions}{Standing assumptions}
The metatheory is ZFC; the Axiom of Choice is used through
\cref{thm:downward-ls}. The set \(\mathrm{Prop}\) of propositional
letters is the countably infinite set of
\cref{def:p03-topological-carrier}, so \(L_2\) and \(L_t\) have countable
signatures. Formulas are finitary. Topological models
\(M=(X,\tau,\nu)\) have \(X\ne\varnothing\), and \(L_t\) is interpreted
on them with open quantifiers ranging over all of \(\tau\). A topological
model is also read as a two-sorted first-order structure: point sort
\(X\), open sort \(\tau\), membership \(\varepsilon\) interpreted by
\(\in\), and \(P_p\) interpreted by \(\nu(p)\). Its \(L_2\)-satisfaction
is exactly first-order satisfaction in that structure. Write
\(\nu_\varnothing\) for the valuation with \(\nu_\varnothing(p)=\varnothing\)
for every \(p\).
\end{hbconvention}

\begin{hbdefinition}{def:p03-w03-p03-051-expressible}{Connectedness and \(L_t\)-expressibility}
A topological space \((X,\tau)\) is \emph{connected} if \(\varnothing\) and
\(X\) are its only subsets that are both open and closed
\cite[Definition~13, p.~11]{TenCateGabelaiaSustretov2006TopologyDraft}.
It is \emph{compact} if every cover of \(X\) by members of \(\tau\) has a
finite subcover.

Let \(K\) be a class of topological models and \(\mathsf P\) a property of
topological spaces. An \(L_t\)-sentence \(\chi\) \emph{expresses
\(\mathsf P\) relative to \(K\)} if, for every \(M=(X,\tau,\nu)\in K\),
\[
 M\models\chi\quad\Longleftrightarrow\quad(X,\tau)\text{ has }\mathsf P.
\]
A set \(\Gamma\) of \(L_t\)-sentences \emph{axiomatizes \(\mathsf P\)
relative to \(K\)} if \(M\models\Gamma\) holds exactly for the members of
\(K\) whose space has \(\mathsf P\). Without qualification \(K\) is the
class of all topological models. Two topological models are
\emph{\(L_t\)-equivalent} if they satisfy the same \(L_t\)-sentences.
\end{hbdefinition}

An expressing sentence \(\chi\) is a one-element axiomatizing set, and
\(\{\neg\chi\}\) then axiomatizes the failure of \(\mathsf P\). So it
suffices to exhibit one model with \(\mathsf P\) and one without it that are
\(L_t\)-equivalent: neither \(\mathsf P\) nor its failure can then be
axiomatized.

\section{Countable elementary pieces of a topological model}
\label{sec:p03-w03-p03-051-pieces}

\begin{hblemma}{thm:p03-w03-p03-051-two-sorted-ls}{proved=annex:p03-w03-p03-051-two-sorted-ls}{Two-sorted downward L\"owenheim--Skolem with parameters}
Assume \cref{def:p03-w03-p03-051-assumptions}. Let \(M=(X,\tau,\nu)\) be a
topological model and let \(A_{\mathrm{pt}}\subseteq X\) and
\(A_{\mathrm{op}}\subseteq\tau\) be countable. There are countable
nonempty sets \(X_0\) and \(O_0\) with
\(A_{\mathrm{pt}}\subseteq X_0\subseteq X\) and
\(A_{\mathrm{op}}\subseteq O_0\subseteq\tau\) such that the two-sorted
structure
\[
 M\restriction(X_0,O_0)=
 \bigl(X_0,O_0,\ \in\cap(X_0\times O_0),\ (\nu(p)\cap X_0)_{p\in\mathrm{Prop}}\bigr)
\]
is an elementary substructure of \(M\): for every \(L_2\)-formula
\(\varphi\), all points \(\bar a\) from \(X_0\) and all opens \(\bar u\)
from \(O_0\),
\[
 M\models\varphi[\bar a,\bar u]
 \quad\Longleftrightarrow\quad
 M\restriction(X_0,O_0)\models\varphi[\bar a,\bar u].
\]
\end{hblemma}

The open sort of \(M\restriction(X_0,O_0)\) consists of subsets of \(X\),
not of \(X_0\). The next lemma replaces each by its trace on \(X_0\).

\begin{hblemma}{thm:p03-w03-p03-051-extent}{proved=annex:p03-w03-p03-051-extent}{The trace model of an elementary piece}
In the situation of \cref{thm:p03-w03-p03-051-two-sorted-ls} put
\[
 e(u)=u\cap X_0\ (u\in O_0),\qquad
 \sigma_0=\{e(u):u\in O_0\},\qquad
 \nu_0(p)=\nu(p)\cap X_0,
\]
and let \(\tau_0\) be the topology on \(X_0\) generated by \(\sigma_0\).
Then:
\begin{enumerate}
\item \(\sigma_0\) is a basis in the sense of
      \cref{def:p03-topological-carrier}, and \(e\) is injective;
\item for every \(L_t\)-formula \(\varphi\), all \(\bar a\) from \(X_0\)
      and all \(\bar u\) from \(O_0\),
      \[
       (X,\tau,\nu)\models\varphi[\bar a,\bar u]
       \quad\Longleftrightarrow\quad
       (X_0,\tau_0,\nu_0)\models\varphi[\bar a,e(\bar u)];
      \]
      in particular the two topological models are \(L_t\)-equivalent;
\item if \(A_{\mathrm{op}}\) is a basis of \(\tau\), then \(\tau_0\) is the
      subspace topology \(\{W\cap X_0:W\in\tau\}\).
\end{enumerate}
\end{hblemma}

Part~(2) is a sharpening of \cref{thm:p03-topological-compactness-ls}: the
countable model is a subspace of the given one, it contains prescribed
points, and it agrees with the given model on formulas with parameters.
Part~(3) is what makes the countable model a recognizable space.

\begin{hbmethod}{def:p03-w03-p03-051-method}{Showing that a property is not \(L_t\)-expressible}
Let \(\mathsf P\) be a property of topological spaces.
\begin{enumerate}
\item Choose a topological model \(M=(X,\tau,\nu)\) whose space has
      \(\mathsf P\) and has a countable basis \(A_{\mathrm{op}}\).
\item Choose a countable \(A_{\mathrm{pt}}\subseteq X\) such that
      \emph{every} countable subspace \(Y\) with
      \(A_{\mathrm{pt}}\subseteq Y\subseteq X\) fails \(\mathsf P\) in its
      subspace topology.
\item Apply \cref{thm:p03-w03-p03-051-two-sorted-ls} and
      \cref{thm:p03-w03-p03-051-extent}: the subspace
      \((X_0,\tau_0,\nu_0)\) is \(L_t\)-equivalent to \(M\) and fails
      \(\mathsf P\).
\end{enumerate}
The same steps apply with the roles of \(\mathsf P\) and its failure
exchanged. Side conditions: the signature must be countable, or the piece
need not be countable; the basis in step~1 must be countable, since it is
placed inside \(O_0\); and step~2 must hold for every countable
\(Y\supseteq A_{\mathrm{pt}}\), because the method gives no further
control over which points \(X_0\) contains.
\end{hbmethod}

\section{The real line and its countable pieces}
\label{sec:p03-w03-p03-051-real-line}

Let \(\tau_{\mathbb R}\) be the usual topology of \(\mathbb R\), and let
\(\mathcal I_{\mathbb Q}=\{\varnothing\}\cup
\{(a,b):a,b\in\mathbb Q,\ a<b\}\) be the rational interval basis,
including the empty set as required by \cref{def:p03-topological-carrier}. For \(Y\subseteq\mathbb R\) write
\(\tau_{\mathbb R}\restriction Y\) for the subspace topology on \(Y\).

\begin{hbtheorem}{thm:p03-w03-p03-051-real-line}{proved=annex:p03-w03-p03-051-real-line}{A disconnected countable subspace \(L_t\)-equivalent to the line}
Assume \cref{def:p03-w03-p03-051-assumptions}. There is a countable set
\(X_0\) with \(\mathbb Q\subseteq X_0\subseteq\mathbb R\) such that:
\begin{enumerate}
\item for every \(L_t\)-formula \(\varphi\), all \(\bar a\) from \(X_0\)
      and all open parameters \(\bar u\) from
      \(\mathcal I_{\mathbb Q}\),
      \[
       (\mathbb R,\tau_{\mathbb R},\nu_\varnothing)\models\varphi[\bar a,\bar u]
       \quad\Longleftrightarrow\quad
       (X_0,\tau_{\mathbb R}\restriction X_0,\nu_\varnothing)
         \models\varphi[\bar a,\bar u\cap X_0];
      \]
\item \((\mathbb R,\tau_{\mathbb R})\) is connected and
      \((X_0,\tau_{\mathbb R}\restriction X_0)\) is not.
\end{enumerate}
\end{hbtheorem}

\begin{hbtheorem}{thm:p03-w03-p03-051-connectedness}{proved=annex:p03-w03-p03-051-connectedness}{Connectedness is not \(L_t\)-expressible}
Assume \cref{def:p03-w03-p03-051-assumptions}. Let \(K_{\mathbb R}\) be the
class of topological models \((Y,\tau_{\mathbb R}\restriction
Y,\nu_\varnothing)\) with \(\varnothing\ne Y\subseteq\mathbb R\).
\begin{enumerate}
\item No set of \(L_t\)-sentences axiomatizes connectedness relative to
      \(K_{\mathbb R}\), and no set axiomatizes its failure.
\item In particular no \(L_t\)-sentence expresses connectedness, relative to
      \(K_{\mathbb R}\) or relative to the class of all topological models.
\item Parts~(1) and~(2) remain true relative to the subclass of
      \(K_{\mathbb R}\) consisting of the models of any set \(\Delta\) of
      \(L_t\)-sentences true in \((\mathbb R,\tau_{\mathbb R},\nu_\varnothing)\).
\end{enumerate}
\end{hbtheorem}

Every space in \(K_{\mathbb R}\) carries the topology of a metric, the
restriction of the usual metric of the line. Part~(3) applies, for
instance, to \(\Delta=\{T_2\}\), the Hausdorff sentence of
\cref{thm:p03-topological-property-formulas}, which is true in the line
(distinct reals \(x<y\) have the disjoint neighborhoods
\((-\infty,m)\) and \((m,\infty)\) with \(m=(x+y)/2\)): adding a
separation axiom expressed in \(L_t\) does not make connectedness
expressible.

\begin{hbcorollary}{thm:p03-w03-p03-051-compactness}{proved=annex:p03-w03-p03-051-compactness}{Compactness is not \(L_t\)-expressible}
Assume \cref{def:p03-w03-p03-051-assumptions}. There is a countable set
\(Z_0\) with \(\mathbb Q\cap[0,1]\subseteq Z_0\subseteq[0,1]\) such that
\(([0,1],\tau_{\mathbb R}\restriction[0,1],\nu_\varnothing)\) and
\((Z_0,\tau_{\mathbb R}\restriction Z_0,\nu_\varnothing)\) are
\(L_t\)-equivalent, the first space is compact, and the second is not.
Hence no set of \(L_t\)-sentences axiomatizes compactness, or its failure,
relative to \(K_{\mathbb R}\); in particular no \(L_t\)-sentence expresses
compactness.
\end{hbcorollary}

\section{A worked example: the unguarded connectedness sentence}
\label{sec:p03-w03-p03-051-example}

Write \(\mathrm{Conn}\) for the \(L_2\)-sentence of
\cref{thm:p03-topological-property-formulas},
\[
 \mathrm{Conn}:\quad
 \forall U\forall V\left[
   \forall x(x\varepsilon U\leftrightarrow\neg(x\varepsilon V))
    \to\bigl(\forall x\,x\varepsilon U\lor
              \forall x\,\neg(x\varepsilon U)\bigr)\right],
\]
which holds in a topological model exactly when its space is connected.
Its open quantifiers are unguarded, and its antecedent uses \(U\) with
both polarities.

\begin{hbcorollary}{thm:p03-w03-p03-051-conn-basis}{proved=annex:p03-w03-p03-051-conn-basis}{The connectedness sentence separates a basis from its topology}
Let \(X_0\), \(O_0\), \(\sigma_0\) be as in the proof of
\cref{thm:p03-w03-p03-051-real-line}. The basoid model
\((X_0,\sigma_0,\nu_\varnothing)\) satisfies \(\mathrm{Conn}\), and the
topological model \((X_0,\tau_{\mathbb R}\restriction X_0,\nu_\varnothing)\)
it generates does not.
\end{hbcorollary}

\begin{hbexample}{def:p03-w03-p03-051-example}{Where the separation happens}
Pick \(r\in\mathbb R\setminus X_0\). In the generated topology the sets
\(L=(-\infty,r)\cap X_0\) and \(R=(r,\infty)\cap X_0\) are complementary
nonempty opens, so they are a counterexample to \(\mathrm{Conn}\). Each is
a union of rational intervals traced on \(X_0\), but by
\cref{thm:p03-w03-p03-051-conn-basis} no two members of the basis
\(\sigma_0\) are complementary in this way. A sentence of \(L_2\) can
therefore see whether the open sort is a basis or the whole topology;
\(\mathrm{Conn}\) is a concrete witness to the last assertion of
\cref{thm:p03-topological-basis}. An \(L_t\)-sentence cannot see the
difference. Conversely, the \(L_t\)-sentence \(T_2\) holds
in \((\mathbb R,\tau_{\mathbb R},\nu_\varnothing)\), and
\cref{thm:p03-w03-p03-051-real-line} transfers it to
\(X_0\), consistent with the fact that \(X_0\) is a metric subspace.
\end{hbexample}

\begin{hbcorollary}{thm:p03-w03-p03-051-modal}{proved=annex:p03-w03-p03-051-modal}{Modal formulas on the countable piece}
For every formula \(\varphi\) of the basic modal language and every
\(a\in X_0\),
\[
 (\mathbb R,\tau_{\mathbb R},\nu_\varnothing),a\models\varphi
 \quad\Longleftrightarrow\quad
 (X_0,\tau_{\mathbb R}\restriction X_0,\nu_\varnothing),a\models\varphi .
\]
Hence no modal formula is true at every point of a model in
\(K_{\mathbb R}\) exactly when that model's space is connected.
\end{hbcorollary}

\begin{hbwarning}{def:p03-w03-p03-051-limits}{What these results do not say}
\begin{enumerate}
\item They concern the finitary guarded language \(L_t\). The unrestricted
      language \(L_2\) expresses connectedness by \(\mathrm{Conn}\). The
      method fails for it at the passage from the piece to its generated
      topology, as \cref{thm:p03-w03-p03-051-conn-basis} shows.
\item The method needs a countable signature and a model with a countable
      basis; with an uncountable signature \cref{thm:downward-ls} does
      not supply a countable piece.
\item The set \(X_0\) is not identified. The proof does not show that
      \(\mathbb Q\) itself, with its subspace topology, is
      \(L_t\)-equivalent to \(\mathbb R\).
\item \cref{thm:p03-w03-p03-051-modal} concerns truth in fixed models with
      the empty valuation. Validity of a modal formula on a space quantifies
      over all valuations, which is second-order over subsets of the
      carrier, and is not treated here.
\item The draft also asserts a converse of basis invariance: every
      \(L_2\)-formula with the invariance property is equivalent on
      topological models to an \(L_t\)-formula
      \cite[Theorem~2.3, p.~9]{TenCateGabelaiaSustretov2006TopologyDraft}.
      The draft does not prove it, it is not proved in this handbook, and
      nothing above uses it.
\end{enumerate}
\end{hbwarning}

\section{Relationships and sources}
\label{sec:p03-w03-p03-051-relationships}

The positive side of the picture is
\cref{thm:p03-topological-property-formulas}: \(T_0\), \(T_1\), \(T_2\),
discreteness and the Alexandroff property are \(L_t\)-expressible, and the
draft's table adds regularity and \(T_3\)
\cite[Table~1, p.~12]{TenCateGabelaiaSustretov2006TopologyDraft}. Each of
these is a local condition, asserting the existence of neighborhoods of
given points with a first-order property. Connectedness and compactness are
global conditions on families of open sets, and the countable pieces of
\cref{thm:p03-w03-p03-051-extent} destroy them while preserving every
\(L_t\)-sentence, including the displayed separation and neighborhood
conditions. The modal corollary uses the standard translation of
\cref{def:p03-topological-modal-translation} and
\cref{thm:p03-topological-translation-homeomorphism}. The ultraproduct
transfer of \cref{thm:p03-topological-transfer} gives a second way to
compare topological models in \(L_t\); it is not needed here.

Sources: \cite[Definition~13, p.~11; Theorem~2.9 and Table~1,
p.~12]{TenCateGabelaiaSustretov2006TopologyDraft} for the definition, the
statement for connectedness and the expressible properties; the draft's
attribution of Theorem~2.9 is to \cite{FlumZiegler1980TopologicalModelTheory}.
